\documentclass[11pt]{article}

%------------------------- fonts & encoding -----------------------------------
\usepackage[T1]{fontenc}
\usepackage[utf8]{inputenc}
\usepackage{amsmath,amssymb,amsthm}
\usepackage{mathtools}
\usepackage[sc]{mathpazo}
\linespread{1.06}
\usepackage{microtype}

%------------------------- layout ---------------------------------------------
\usepackage[a4paper,margin=1.05in,marginparwidth=1.0in]{geometry}
\usepackage{booktabs}
\usepackage{enumitem}
\setlist{itemsep=2pt,topsep=3pt}
\usepackage{xcolor}
\usepackage{graphicx}
\usepackage{tikz}
\usepackage{pgfplots}
\pgfplotsset{compat=1.17}
\usepackage[most,breakable]{tcolorbox}

%------------------------- palette --------------------------------------------
\definecolor{crimson}{HTML}{A51C30}
\definecolor{ink}{HTML}{1A1A1A}
\definecolor{steel}{HTML}{2C5F8A}
\definecolor{amber}{HTML}{9A6A00}
\definecolor{teal}{HTML}{1F6F6B}
\definecolor{paperlt}{HTML}{F7F4EF}

%------------------------- section styling ------------------------------------
\usepackage{titlesec}
\titleformat{\section}{\Large\bfseries\color{steel}}{\thesection}{0.7em}{}
\titleformat{\subsection}{\large\bfseries\color{ink}}{\thesubsection}{0.6em}{}
\titlespacing*{\section}{0pt}{1.6em}{0.6em}

%------------------------- hyperref -------------------------------------------
\usepackage[colorlinks=true,linkcolor=steel,citecolor=steel,urlcolor=crimson]{hyperref}

%------------------------- theorem environments -------------------------------
\theoremstyle{definition}
\newtheorem{definition}{Definition}[section]
\newtheorem{example}{Example}[section]
\newtheorem{exercise}{Exercise}[section]
\theoremstyle{plain}
\newtheorem{proposition}{Proposition}[section]
\theoremstyle{remark}
\newtheorem*{remark}{Remark}

%------------------------- callout boxes --------------------------------------
\newtcolorbox{intuition}{
  enhanced, breakable, colback=crimson!5, colframe=crimson, boxrule=0.5pt,
  left=8pt,right=8pt,top=6pt,bottom=6pt, arc=2pt,
  colbacktitle=crimson, coltitle=paperlt, fonttitle=\bfseries, title={Intuition}}

\newtcolorbox{pitfall}{
  enhanced, breakable, colback=amber!7, colframe=amber, boxrule=0.5pt,
  left=8pt,right=8pt,top=6pt,bottom=6pt, arc=2pt,
  colbacktitle=amber, coltitle=paperlt, fonttitle=\bfseries, title={Pitfall}}

\newtcolorbox{keyresult}{
  enhanced, breakable, colback=steel!6, colframe=steel, boxrule=0.5pt,
  left=8pt,right=8pt,top=6pt,bottom=6pt, arc=2pt,
  colbacktitle=steel, coltitle=paperlt, fonttitle=\bfseries, title={Takeaway}}

\newtcolorbox{observed}{
  enhanced, breakable, colback=teal!6, colframe=teal, boxrule=0.5pt,
  left=8pt,right=8pt,top=6pt,bottom=6pt, arc=2pt,
  colbacktitle=teal, coltitle=paperlt, fonttitle=\bfseries, title={Measured}}

%------------------------- math operators -------------------------------------
\DeclareMathOperator{\Var}{Var}
\DeclareMathOperator{\Cov}{Cov}
\DeclareMathOperator{\Corr}{corr}
\DeclareMathOperator{\diag}{diag}
\DeclareMathOperator{\std}{std}
\DeclareMathOperator*{\avg}{\mathbb{E}}
\newcommand{\T}{^{\!\top}}
\newcommand{\R}{\mathbb{R}}
\newcommand{\half}{\tfrac12}
\newcommand{\wvec}{\mathbf{w}}
\newcommand{\Xmat}{\mathbf{X}}
\newcommand{\Fmat}{\mathbf{F}}
\newcommand{\Vmat}{\mathbf{V}}
\newcommand{\Umat}{\mathbf{U}}
\newcommand{\Rmat}{\mathbf{R}}
\newcommand{\Lam}{\boldsymbol{\Lambda}}
\newcommand{\Del}{\boldsymbol{\Delta}}
\newcommand{\fvec}{\mathbf{f}}

%------------------------- title ----------------------------------------------
\title{\vspace{-2.2em}
  {\color{steel}\rule{\textwidth}{1.3pt}}\\[0.7em]
  \Huge\bfseries Growth\\[0.35em]
  \large\mdseries The trend factor --- how fast a firm's fundamentals are rising, measured as a normalised slope\\[0.4em]
  {\color{steel}\rule{\textwidth}{0.7pt}}}
\author{}
\date{}

%------------------------- watermark footer -----------------------------------
\usepackage{fancyhdr}
\pagestyle{fancy}
\fancyhf{}
\renewcommand{\headrulewidth}{0pt}
\renewcommand{\footrulewidth}{0pt}
\fancyfoot[L]{\footnotesize\color{black!45}MSCI Barra USE4 Imitation}
\fancyfoot[R]{\footnotesize\color{black!45}\thepage}

\begin{document}
\maketitle
\thispagestyle{empty}

\begin{abstract}
Growth is the trend factor: it measures how fast a firm's fundamentals are rising, so that fast-growing companies load high and shrinking ones load low. USE4 builds Growth from a blend of growth descriptors -- a forward, analyst-predicted earnings-growth view (the heaviest-weighted ingredient), a trailing earnings-growth view, and a trailing sales-growth view -- each computed as a \emph{normalised} multi-year regression slope and then standardised across the estimation universe. Normalising the slope by the average level is the load-bearing idea: it converts a dollars-per-year trend into a scale-free growth \emph{rate}, so a small firm adding a few cents of earnings a year is not ranked against a giant adding dollars. Our personal build cannot source the forward analyst view, so it drops that dominant ingredient and rests the factor on realised, backward-looking growth -- a faithful trend signal, but a realised one rather than an expected one. The construction also exposes the one place growth is genuinely undefined: a firm that loses money on average over the window has no meaningful growth \emph{rate}, because dividing a slope by a non-positive mean is meaningless. The result is a sharply peaked, fat-tailed cross-section of growth rates that the trim and standardisation must tame.
\end{abstract}

\begin{tcolorbox}[colback=paperlt,colframe=ink!30,boxrule=0.4pt,arc=2pt,left=10pt,right=10pt,top=8pt,bottom=8pt]
\textbf{How to read these notes.} We move from \emph{why} Growth is a factor to \emph{how} USE4 specifies it, then to how we actually compute it from the data we have, with a deliberate stop at the normalisation step (Section~\ref{sec:slope}) and at the one case where growth has no meaning (Section~\ref{sec:lossmaker}). Read those two slowly; they carry the weight of the set. The numbers labelled as measured come from a single run of our personal build.

\smallskip
Four coloured boxes recur: \textcolor{crimson}{\textbf{Intuition}} gives the mental picture, \textcolor{amber}{\textbf{Pitfall}} flags what breaks without the fix, \textcolor{steel}{\textbf{Takeaway}} states a load-bearing result, and \textcolor{teal}{\textbf{Measured}} reports real numbers from the build run.

\smallskip
\textbf{Reference material.}
\begin{itemize}[leftmargin=1.4em,topsep=2pt]
  \item External: Menchero, Orr \& Wang (2011), USE4 Empirical Notes (Appendix A) and Methodology Notes (Sec. 2.2--2.3).
\end{itemize}
\end{tcolorbox}

\tableofcontents
\bigskip

\section{What Growth is, and why it is a factor}\label{sec:what}

\paragraph{Learning objectives.}
\begin{itemize}[leftmargin=1.4em,topsep=2pt]
  \item State Growth as the trend in a firm's fundamentals, not their level.
  \item Explain why a \emph{rate} of growth, not a dollar amount of growth, is the right quantity.
  \item Say in one sentence what the Growth exposure is, before any standardisation.
\end{itemize}

Growth is about direction, not position. Where value asks whether a firm is cheap and size asks whether it is large, growth asks whether its business is \emph{getting bigger}, and how quickly. Two firms can be identical in size and valuation today and still be very different bets: one with earnings climbing year after year, the other with earnings flat or fading. That trajectory is a distinct dimension of risk and return -- growth stocks tend to move together, re-rate together, and disappoint together -- and so it earns its own factor.

The natural way to measure a trend is its slope: fit a line to several years of a firm's earnings and read off how fast they are rising. But a raw slope measured in dollars is not comparable across firms. A company earning a few cents a share that adds a penny a year is growing explosively; a company earning ten dollars a share that adds the same penny is barely moving. The dollar slopes are identical; the growth \emph{rates} are worlds apart. So the quantity we want is not the slope itself but the slope \emph{relative to the firm's own scale} -- a normalised slope, which behaves like an average annual growth rate.

\begin{intuition}
The whole factor answers one question: \emph{how fast are this firm's fundamentals trending up, as a fraction of where they sit?} A high number means earnings (and sales) climbing briskly relative to their own level; a number near zero means a flat business; a negative number means a fading one. Hold that picture --- a growth \emph{rate}, not a dollar trend --- and every computational choice below follows.
\end{intuition}

\begin{keyresult}
The Growth exposure is a standardised blend of normalised multi-year growth slopes -- chiefly of earnings, with sales as a supporting view. Everything that follows is about which series, over what window, and how to make the slope scale-free.
\end{keyresult}

\section*{Notation}
\addcontentsline{toc}{section}{Notation}

\begin{center}
\begin{tabular}{@{}ll@{}}
\toprule
Symbol & Meaning \\
\midrule
$e_{i,\tau}$ & annual earnings per share of stock $i$ in window-year $\tau$ \\
$s_{i,\tau}$ & annual sales per share of stock $i$ in window-year $\tau$ \\
$\beta^{E}_i,\ \beta^{S}_i$ & OLS slope of $e_{i,\tau}$ (resp.\ $s_{i,\tau}$) on the year index \\
$\bar e_i,\ \bar s_i$ & mean earnings (resp. sales) per share over the window \\
$\text{EGRO}_i$ & normalised earnings-growth slope, $\beta^{E}_i / \bar e_i$ \\
$\text{SGRO}_i$ & normalised sales-growth slope, $\beta^{S}_i / \bar s_i$ \\
$\text{EGRLF}_i$ & forward analyst earnings-growth forecast (unavailable in our data) \\
$g_i$ & raw GRO descriptor: a weighted blend of $\text{EGRO}_i$ and $\text{SGRO}_i$ \\
$x_i$ & standardised GRO exposure for stock $i$ \\
$w_i$ & cap weight of stock $i$ in the ESTU \\
$\avg_{\text{cw}}, \std_{\text{ew}}$ & cap-weighted mean, equal-weighted standard deviation over the ESTU \\
$Q_1,\dots,Q_5$ & quintiles of raw $g_i$, from most contracting ($Q_1$) to fastest-growing ($Q_5$) \\
ESTU & estimation universe (MSCI USA IMI proxy) \\
\bottomrule
\end{tabular}
\end{center}

\section{How USE4 defines it}\label{sec:def}

\paragraph{Learning objectives.}
\begin{itemize}[leftmargin=1.4em,topsep=2pt]
  \item Name the growth descriptors USE4 blends, and which one it weights most.
  \item State that each descriptor is a normalised multi-year slope, and the standardisation target.
\end{itemize}

USE4 builds Growth as a blend of growth descriptors (Menchero, Orr \& Wang 2011, Empirical Notes, App. A). The dominant ingredient is a \emph{forward} view: analysts' predicted long-term earnings growth (EGRLF), which USE4 weights most heavily, carrying the bulk of the blend. The remainder is split between two \emph{trailing} views -- realised earnings growth (EGRO) and realised sales growth (SGRO) -- with earnings growth the larger of the two. In the published specification the forward view dominates and the trailing earnings view is secondary, with sales growth a small supporting term.

Each trailing descriptor is a \emph{normalised} multi-year regression slope: fit an ordinary least-squares line to several years of annual per-share earnings (or sales), take the slope, and divide by the mean level over the same window. Dividing by the mean is what makes the descriptor a scale-free growth rate rather than a dollars-per-year trend. The blended descriptor is then standardised over the estimation universe to cap-weighted mean $0$ and equal-weighted standard deviation $1$ (Methodology Notes, Sec. 2.3), after a three-standard-deviation trim (Sec. 2.2). Fundamentals are aligned point-in-time by filing date, so any signal date uses only earnings a market participant could actually have known.

\section{How we compute it: the forward hole, and a normalised slope}\label{sec:compute}

\paragraph{Learning objectives.}
\begin{itemize}[leftmargin=1.4em,topsep=2pt]
  \item Explain why the forward view is dropped, that it was the heaviest-weighted view, and what that does to the factor.
  \item Reconstruct the normalised-slope descriptor from annual per-share data.
  \item Identify the computational choices that most change a firm's ranking.
\end{itemize}

The definition assumes three clean inputs. Our data supplies two of them and not the third.

\paragraph{(a) No forward earnings growth.} The data carries no analyst long-term growth forecasts, so the EGRLF view cannot be computed and is dropped. This is the most material deviation in the whole factor, because the dropped view is not a minor third ingredient -- it is the one USE4 weights most heavily, the dominant component of the intended blend. With it gone, the forward weight collapses onto the trailing earnings-growth view, leaving an effective blend of $0.90\,\text{EGRO}_i + 0.10\,\text{SGRO}_i$: nine parts realised earnings growth to one part realised sales growth. The consequence is plain and worth stating out loud: our Growth factor is \emph{realised}, not \emph{expected}. USE4's growth factor is, by weight, mostly a statement about what analysts think a firm will do; ours is entirely a statement about what the firm has already done. The two usually agree for steady compounders and disagree most for turnarounds and young firms, where expectations run ahead of (or behind) the trailing record. There is a subtler consequence too: analysts' long-term forecasts are almost always positive, so the forward view USE4 leans on would pull the blend upward; without it, our realised-only growth sits lower and more negatively tilted than full USE4 would --- consistent with the mostly-negative cross-section reported in Section~\ref{sec:std}.

\paragraph{(b) The normalised slope.} For each stock at each signal date we take the most recent annual per-share earnings and sales -- a window of up to five years, requiring at least a few annual points so a slope is well determined -- using only filings whose report date falls on or before the signal date. The underlying fundamentals are reported quarterly, so each annual earnings or sales figure is itself aggregated from the firm's quarterly filings. We regress each series on the year index, take the OLS slope, and divide by the mean level over the window:
\[
\text{EGRO}_i \;=\; \frac{\beta^{E}_i}{\bar e_i},
\qquad
\text{SGRO}_i \;=\; \frac{\beta^{S}_i}{\bar s_i}.
\]
The slope $\beta^{E}_i$ answers ``how many dollars per year''; dividing by $\bar e_i$ turns it into ``what fraction of the typical year's earnings per year'' -- an average annual growth rate, comparable across a penny-earner and a ten-dollar-earner alike. Each descriptor is then clipped to a fixed band before blending, because a normalised slope can blow up when the mean level sits close to zero (more on that in Section~\ref{sec:lossmaker}); the clip caps those explosions without discarding the names.

\paragraph{(c) The blend, and stale filers.} The two descriptors are combined with the effective $0.90 / 0.10$ weights; if one descriptor is missing, the available weight is renormalised to one so the surviving view still produces a score. Finally, a firm whose most recent filing is too old -- older than roughly a year and a half -- is treated as having no current growth signal, rather than carrying a long-dead filing forward. This staleness rule matters more than it sounds: without it, firms that stopped reporting years ago keep a frozen growth score indefinitely, quietly polluting the cross-section with stale data.

\begin{pitfall}
Three choices here move rankings more than they look. \emph{Forgetting to normalise} the slope ranks firms by dollars-per-year, which is almost a restatement of size -- big firms swamp the cross-section and the growth signal is lost. \emph{Skipping the clip} lets a near-zero mean detonate the ratio, handing a single borderline firm an absurd score. And \emph{ignoring staleness} lets dead filers keep a fossil growth score forever. The quietest trap, though, is the one you cannot see in the output: nothing flags that the forward view is missing, so the renormalisation silently rests the whole factor on backward-looking data -- correct, but invisible unless stated out loud.
\end{pitfall}

\section{Why normalise by the mean: the slope alone is not growth}\label{sec:slope}

\paragraph{Learning objectives.}
\begin{itemize}[leftmargin=1.4em,topsep=2pt]
  \item Show that a raw slope is not scale-free, and that dividing by the mean fixes it.
  \item Explain what the normalised slope approximates (an average annual growth rate).
\end{itemize}

This is the step that turns a regression into a factor. A slope measured in dollars carries the firm's scale baked in, so ranking by raw slope ranks mostly by size. The normalisation strips the scale out. Concretely, if earnings grow at a roughly constant fractional rate, the slope is proportional to the level, and dividing the slope by the mean level cancels the scale, leaving the rate itself. That is why $\beta^{E}_i / \bar e_i$ is, to a good approximation, the firm's average annual earnings growth rate over the window -- a pure number you can compare across the whole market.

\begin{example}
Firm A earns \$0.20 a share and has been adding \$0.04 a year; firm B earns \$10.00 a share and has been adding the same \$0.04 a year. Their raw slopes are identical at \$0.04/year. But $\text{EGRO}_A = 0.04/0.20 = 0.20$ -- twenty percent annual growth -- while $\text{EGRO}_B = 0.04/10.00 = 0.004$, four-tenths of a percent. The normalisation correctly calls A a fast grower and B nearly flat; the raw slope would have called them equal.
\end{example}

\begin{intuition}
Think of the normalised slope as the answer to ``if this trend continued, by roughly what percentage would earnings grow each year?'' That framing also explains why the descriptor is fragile near zero earnings: a percentage growth rate needs a sensible base, and a firm whose average earnings are near zero has no sensible base to grow \emph{from}. We meet that failure head-on next.
\end{intuition}

\section{When growth has no meaning: loss-makers}\label{sec:lossmaker}

\paragraph{Learning objectives.}
\begin{itemize}[leftmargin=1.4em,topsep=2pt]
  \item Explain why a normalised earnings-growth slope is undefined when average earnings are non-positive.
  \item Justify setting the descriptor to missing for those names rather than forcing a number.
  \item Contrast this with how an earnings-\emph{yield} factor treats the same loss-makers.
\end{itemize}

The normalisation that makes the slope a growth \emph{rate} also breaks it for an important population. If a firm's average earnings over the window are zero or negative -- a chronic or deep loss-maker -- then dividing the slope by that mean is meaningless. The sign flips perversely (a loss that is shrinking, an improvement, divided by a negative mean, reads as negative ``growth''), and the magnitude explodes as the mean approaches zero. A growth \emph{rate} simply does not exist when there is no positive base to grow from.

The build's response is principled: when the average earnings level over the window is non-positive, the earnings-growth descriptor is set to missing rather than forced into a number. If sales growth is still available -- sales are almost always positive even for an unprofitable firm -- the blend renormalises onto it; if neither view is computable, the firm has no growth score that month. The principle is that a fabricated growth rate is worse than an absent one: a nonsensical ratio would inject pure noise into exactly the names whose fundamentals are most unstable.

\begin{pitfall}
Do not ``rescue'' a loss-maker's growth rate by dividing the slope by the \emph{absolute value} of a negative mean, or by flooring the mean at a small positive number. Both manufacture a finite figure where no growth rate exists, and both put their largest fabricated values on the least stable firms. The correct move is to mark the descriptor missing and let the blend lean on sales -- or let the name go unscored -- not to invent a rate.
\end{pitfall}

\begin{intuition}
This is the mirror image of how an earnings-\emph{yield} factor treats the same firms. Earnings yield $E/P$ \emph{keeps} loss-makers, because a negative yield is real, informative, and well defined -- a price divided by a loss is simply a negative number. Growth \emph{cannot} keep them, because a growth \emph{rate} needs a positive base and a loss has none. Same firms, opposite handling, and for a precise reason: a ratio to price stays meaningful through zero earnings; a ratio to average earnings does not.
\end{intuition}

\section{Trim and standardisation}\label{sec:std}

\paragraph{Learning objectives.}
\begin{itemize}[leftmargin=1.4em,topsep=2pt]
  \item Apply the three-sigma trim and explain why growth slopes stay fat-tailed.
  \item State the standardisation target and read the stability diagnostics.
\end{itemize}

The raw blended growth rate $g_i$ is trimmed at three standard deviations and then standardised over the ESTU to cap-weighted mean $0$ and equal-weighted standard deviation $1$:
\[
x_i \;=\; \frac{g_i - \avg_{\text{cw}}(g)}{\std_{\text{ew}}(g)},
\qquad
\avg_{\text{cw}}(g) = \frac{\sum_i w_i\, g_i}{\sum_i w_i}.
\]
The cap-weighted centring means a stock's exposure is measured against the dollar-weighted market's growth rate; the equal-weighted scaling fixes the spread. Growth is, even after the per-descriptor clip, an unusually fat-tailed quantity: a handful of firms genuinely grow (or collapse) several times faster than the typical name, and a normalised slope amplifies anything near a small base. So the three-sigma trim does more work here than for a tame descriptor -- and still leaves heavier tails than usual.

\begin{observed}
On the build run the raw growth cross-section is sharply peaked at a slightly negative centre -- the median normalised slope is $\approx -0.08$, and about two-thirds of in-universe names sit below zero (the typical firm's per-share earnings are not, on a normalised five-year slope, trending up). The distribution is very fat-tailed: excess kurtosis $\approx 8$, with a mild left skew. The three-sigma trim clips $\approx 2.75\%$ of ESTU rows -- above the usual $0.5$--$2\%$ band, a direct symptom of those fat tails. After standardisation the cap-weighted mean is $0$ to machine precision and the equal-weighted standard deviation is $1.0$, with $\approx 96\%$ of ESTU rows inside $[\pm 3]$. The exposure is computed for every stock but standardised only on the ESTU, so that in-universe fraction is a touch higher than the all-stocks figure ($\approx 95.6\%$): out-of-universe names are scored but do not calibrate the mean and spread. Sorting into quintiles by raw growth gives a strictly monotone profile, from a contracting bottom quintile (mean slope $\approx -0.8$) up to a fast-growing top ($\approx +0.6$), a spread of $\approx 1.41$. Month-over-month rank stability is high, Spearman $\rho \approx 0.97$: a five-year slope moves slowly as one new annual point replaces an old one.
\end{observed}

\begin{figure}[htbp]
\centering
\includegraphics[width=\linewidth]{figures/gro_dist.png}
\caption{The raw growth cross-section -- normalised five-year slopes of per-share earnings and sales -- pooled over in-ESTU stock-dates. The distribution is sharply peaked at a slightly negative median ($\approx -0.08$) with conspicuously fat tails: a few firms grow or contract several times faster than the typical name. The shaded regions mark a representative three-sigma trim band; growth stays fat-tailed enough that the trim clips $\approx 2.75\%$ of ESTU rows, above the usual $0.5$--$2\%$.}
\label{fig:gro_dist}
\end{figure}

\begin{observed}
The deliverable spans $\approx 1.46$M stock-month rows over $\approx 315$ monthly dates (it begins more than a year later than the price-based factors, since a multi-year slope needs several annual filings to exist) and $\approx 13{,}450$ tickers; the raw descriptor is missing for only $\approx 2.6\%$ of rows, and most scored names use the full five-year window.
\end{observed}

\begin{exercise}
A firm's annual EPS over five years is \$1.00, \$1.10, \$1.20, \$1.30, \$1.40. Estimate the OLS slope (the yearly increment is constant here) and the mean, then form EGRO. Roughly what annual growth rate does it imply, and does that match your intuition for a series rising \$0.10 a year off a \$1.20 average?
\end{exercise}

\begin{exercise}
Two firms have identical raw earnings slopes of \$0.50/year. Firm A averaged \$2.00 of EPS over the window; firm B averaged \$25.00. Which loads higher on Growth, and by how much? State in one sentence what would happen to the cross-section if the build ranked by raw slope instead of normalised slope.
\end{exercise}

\begin{exercise}
Spearman $\rho \approx 0.97$ month-over-month for Growth. Explain why a five-year normalised slope should be \emph{more} month-to-month stable than a single-quarter earnings number, but \emph{less} stable than a factor built purely from market cap. (Hint: ask how much new information one extra annual filing adds to a five-point regression.)
\end{exercise}

\section{Summary, and one caveat}\label{sec:summary}

\paragraph{Learning objectives.}
\begin{itemize}[leftmargin=1.4em,topsep=2pt]
  \item Recite the Growth build from annual filings to standardised exposure.
  \item State the central limitation and the lever that would address it.
\end{itemize}

The build, end to end: gather up to five years of annual per-share earnings and sales, point-in-time; regress each on the year index and divide the slope by its mean level to form a scale-free growth rate; clip each descriptor to tame near-zero-base explosions; set the earnings descriptor missing where average earnings are non-positive, since a growth rate needs a positive base; blend nine parts earnings growth to one part sales growth, renormalising if one view is absent, and dropping names whose newest filing is too stale; trim at three standard deviations; standardise to cap-weighted mean $0$, equal-weighted standard deviation $1$. The output is a sharply peaked, fat-tailed, slowly-moving trend factor.

\begin{keyresult}
Growth is the normalised trend in fundamentals. The two computations that matter most are the normalisation -- dividing the slope by the mean level, which turns a dollar trend into a comparable growth rate -- and the loss-maker rule, which refuses to manufacture a growth rate where no positive base exists.
\end{keyresult}

\begin{pitfall}
One honest limitation dominates. This Growth factor is purely \emph{realised}: the analyst-forecast view that USE4 weights most heavily has been dropped, so the factor measures what a firm has already done, not what it is expected to do. It will therefore mark a genuine turnaround as a slow grower until the realised record catches up, and it lacks the forward signal that would have separated firms the market expects to accelerate from those it expects to stall. The single lever that would close most of this gap is a source of analyst long-term growth estimates, which would restore the dominant forward descriptor and make the factor expectational rather than purely backward-looking.
\end{pitfall}

\begin{remark}
None of this changes the factor's place in the model: Growth remains the trend dimension, sitting alongside value, size, and the rest. The caveat is about \emph{timeliness} -- realised versus expected growth -- not about whether the trend signal is real.
\end{remark}

\begin{center}\small\itshape
References: Menchero, Orr \& Wang (2011), \textit{The Barra US Equity Model (USE4)} --- Empirical Notes (Appendix A) and Methodology Notes (Sec. 2.2--2.3). ESTU treated as an MSCI USA IMI proxy; measured quantities are from a personal build run.
\end{center}

\end{document}
